\documentclass[spanish,letterpaper,oneside]{article}

\def\documenttitle {Segundo Momento: Cuadro Diagnóstico de la Práctica Docente Matemática}
\def\documentsubtitle {Fundamentos para la enseñanza y el aprendizaje I · Subproducto 6B (Bloque III)}
\def\documentsubject {Subproducto Independiente 6B · Segundo Momento (Guía 51705 / 51782)}
\def\documentauthor {Mtro. Martín Jonathan de la Cruz Muñoz}
\def\coursename {Fundamentos para la enseñanza y el aprendizaje I}
\def\coursecode {MAEM-SD26}

\def\universityname {Instituto de Investigación, Innovación y Estudios de Posgrado para la Educación del Estado de Nuevo León}
\def\universityfaculty {Maestría en Aprendizaje y Enseñanza de las Matemáticas}
\def\universitydepartment {Fundamentos para la enseñanza y el aprendizaje I}
\def\universitydepartmentimage {img/iiiepe/marca-iiiepe.png}
\def\universitydepartmentimagecfg {height=1.75cm}
\def\universitylocation {Monterrey, Nuevo León}
\def\studentfullname {Mtro. Martín Jonathan de la Cruz Muñoz}
\def\studentid {00874}
\def\studentemail {martin.delacruz@campus.iiiepe.edu.mx}
\def\programperiod {2do periodo, septiembre--diciembre de 2026 (SD26)}
\def\programmodality {Modalidad presencial}
\def\programshort {MAEM}

\def\authortable {
  \begin{tabular}{ll}
    \textbf{Alumno:} & \studentfullname \\
    Matrícula: & \studentid \\
    Correo institucional: & \studentemail \\
    Unidad: & \coursename \\
    Programa: & \programshort \\
    Periodo: & \programperiod \\
    & \\
    \multicolumn{2}{l}{Fecha de realización: 19 de noviembre de 2026} \\
    \multicolumn{2}{l}{\universitylocation}
  \end{tabular}
}

\input{template}

\usepackage{helvet}
\renewcommand{\familydefault}{\sfdefault}
\usepackage{array}
\usepackage{longtable}
\usepackage{pdflscape}
\usepackage{tikz}
\usepackage{pgfplots}
\pgfplotsset{compat=1.18}
\usepackage[most]{tcolorbox}
\usetikzlibrary{arrows.meta,positioning,calc,fit,backgrounds,shapes.geometric}

\definecolor{iiiepeNavy}{HTML}{163B5C}
\definecolor{iiiepeBlue}{HTML}{1F6F9F}
\definecolor{iiiepeAqua}{HTML}{20B7A8}
\definecolor{iiiepeOrange}{HTML}{F28C28}
\definecolor{iiiepePaper}{HTML}{F6F8F7}
\definecolor{iiiepeInk}{HTML}{1E2A32}
\definecolor{iiiepePale}{HTML}{EEF5F8}
\definecolor{iiiepeMint}{HTML}{E7F6F4}
\definecolor{iiiepeWarm}{HTML}{FFF3E6}
\def\portraittitlecolor {iiiepeNavy}
\setcitestyle{authoryear,open={(},close={)}}
\makeatletter
\renewcommand{\@biblabel}[1]{}
\makeatother
\setlength{\bibhang}{1.27cm}
\setlength{\bibsep}{6pt}
\AtBeginDocument{\renewcommand{\refname}{Referencias}}

\newcommand{\insertcoverwatermark}{%
  \AddToShipoutPictureBG*{%
    \begin{tikzpicture}[remember picture,overlay]
      \node[opacity=0.04,inner sep=0pt] at (current page.center) {%
        \includegraphics[width=0.65\paperwidth]{img/iiiepe/marca-iiiepe.png}%
      };
    \end{tikzpicture}%
  }%
}

\setlength{\parskip}{5pt}\setlength{\parindent}{0pt}
\begin{document}\insertcoverwatermark\templatePortrait\templatePagecfg\templateFinalcfg\fancyfoot{}\fancyfoot[L]{\scriptsize IIIEPE · MAEM · SD26}\fancyfoot[C]{\scriptsize\thepage}\fancyfoot[R]{\scriptsize Subproducto 6B · Segundo Momento}\setstretch{1.3}\raggedbottom
\section{Segundo momento: diagnóstico honesto de la práctica de enseñanza}
De conformidad con la segunda etapa de la \textit{Guía instruccional del Bloque III} (módulos \texttt{51705} / \texttt{51782}), presento de forma autónoma el análisis diagnóstico de mi propia práctica docente en educación media superior técnica (CONALEP y CECyTE) y formación universitaria en ingeniería.

Mi postura es que todo rediseño pedagógico honesto requiere identificar primero las propias inercias metodológicas y evaluativas mediante evidencias reales observadas en el aula, en lugar de atribuir las dificultades de aprendizaje exclusivamente a carencias previas del estudiantado \citep{quiroga2022aprendizaje,restrepo2004investigacion}. En consecuencia, el foco problemático elegido se sitúa en el aprendizaje de la \textbf{modelación con funciones lineales y diferenciación entre relación proporcional directa ($y=kx$) y función afín ($y=mx+b$)}.

\section{Cuadro diagnóstico de los seis aspectos críticos}
Al revisar cuadernos de clase, guías de ejercicios y listas de cotejo de secuencias previas, observo un patrón nítido: en 18 de aproximadamente 25 cuadernos revisados, los alumnos ejecutan correctamente sustituciones algebraicas de fórmula cerrada cuando conocen los pasos, pero fracasan al explicar qué significan físicamente la pendiente y la ordenada al origen o al decidir cuándo un modelo deja de ser proporcional. Además, el cuadro siguiente sistematiza este hallazgo a través de las seis categorías de autodiagnóstico docente:

{\small\setstretch{1.15}\renewcommand{\arraystretch}{1.35}
\begin{longtable}{|>{\raggedright\arraybackslash}p{0.20\linewidth}|>{\raggedright\arraybackslash}p{0.25\linewidth}|>{\raggedright\arraybackslash}p{0.49\linewidth}|}\hline
\rowcolor{iiiepePale} \textbf{Aspecto diagnóstico} & \textbf{Pregunta orientadora} & \textbf{Análisis situado con evidencias concretas} \\\hline
\endfirsthead
\hline
\rowcolor{iiiepePale} \textbf{Aspecto diagnóstico} & \textbf{Pregunta orientadora} & \textbf{Análisis situado con evidencias concretas} \\\hline
\endhead
1. Problemática identificada & ¿Qué dificultades presentan los alumnos al aprender matemáticas? & En 18 de aproximadamente 25 cuadernos revisados sobre rectas lineales, el estudiantado obtiene $m$ y $b$ mecánicamente por fórmula, pero no interpreta su sentido físico ni distingue $y=kx$ frente a $y=mx+b$ \citep{star2015algebra}. \\\hline
2. Práctica de enseñanza actual & ¿Cómo abordo la enseñanza de ese contenido matemático? & Suelo iniciar definiendo $y=mx+b$ en pizarrón, resuelvo dos casos modelo de rutina y asigno series individuales antes de proponer aplicaciones en taller o laboratorio. \\\hline
3. Concepciones pedagógicas & ¿Qué ideas sobre enseñar y aprender matemáticas orientan mi práctica? & Subyace la creencia ingenieril de que dominar primero el algoritmo descontextualizado prepara para aplicar después, olvidando que el significado físico construye el concepto algebraico \citep{sawyer2008optimizing}. \\\hline
4. Características del aprendizaje & ¿Qué evidencias muestran aprendizaje superficial o profundo? & Superficial: resolución correcta en formatos rígidos pero parálisis ante cambio de variables o pregunta interpretativa. Profundo incipiente: conjeturas espontáneas en taller cuando vinculan pendiente con sensibilidad real. \\\hline
5. Causas posibles & ¿Qué decisiones didácticas contribuyen a la problemática? & Priorizar ejercitación sin contexto, resolver dudas anticipadamente por presión del reloj escolar y ponderar casi solo exactitud numérica final \citep{eef2017mathematics}. \\\hline
6. Necesidad de transformación & ¿Qué tendría que modificar para favorecer aprendizaje profundo? & Invertir la secuencia: iniciar con un reto técnico con datos, activar colaboración con roles rotativos, ofrecer andamiaje gradual y retroalimentar antes de calificar \citep{eef2021feedback}. \\\hline
\end{longtable}}

\clearpage
\section{Conclusiones}
En conclusión, sostengo que este Segundo momento transforma una queja genérica («el alumnado no domina el álgebra») en un diagnóstico docente accionable y verificable. Considero que reconocer cómo mi propia rapidez expositiva y las tareas cerradas fomentaron aprendizaje superficial brinda la base ética e intelectual necesaria para diseñar la propuesta de innovación del Tercer momento \citep{restrepo2004investigacion,quiroga2022aprendizaje}.

\clearpage
\begin{thebibliography}{99}
\bibitem[Collin y Quigley(2021)]{eef2021feedback} Collin, J., y Quigley, A. (2021). \textit{Teacher feedback to improve pupil learning}. Education Endowment Foundation.
\bibitem[Henderson et~al.(2017)]{eef2017mathematics} Henderson, P., Hodgen, J., Foster, C., y Kuchemann, D. (2017). \textit{Improving mathematics in Key Stages 2 and 3}. Education Endowment Foundation.
\bibitem[Quiroga Aguilar y Lara Órdenes(2022)]{quiroga2022aprendizaje} Quiroga Aguilar, L., y Lara Órdenes, E. (2022). El aprendizaje profundo como herramienta para cambio en la visión de aprendizaje de una cultura escolar. \textit{Revista Educación las Américas}, 12(1).
\bibitem[Restrepo Gómez(2004)]{restrepo2004investigacion} Restrepo Gómez, B. (2004). La investigación-acción educativa y la construcción de saber pedagógico. \textit{Educación y Educadores}, (7), 45--55.
\bibitem[Sawyer(2008)]{sawyer2008optimizing} Sawyer, R. K. (2008). Optimizing learning: Implications of learning sciences research. OECD Publishing.
\bibitem[Star et~al.(2015)]{star2015algebra} Star, J. R., et al. (2015). \textit{Teaching strategies for improving algebra knowledge}. IES.
\end{thebibliography}

\end{document}
